\documentclass[a4paper,12pt]{article}
\usepackage[utf8]{inputenc}
\usepackage[T2A]{fontenc}
\usepackage[russian]{babel}
\usepackage{amsmath}
\usepackage{amssymb}
\usepackage{array}
\usepackage{booktabs}
\usepackage{geometry}
\usepackage{xcolor}
\usepackage{tabularx}
\usepackage{multirow}
\usepackage{makecell}

\geometry{top=2cm, bottom=2cm, left=2.5cm, right=2cm}

\title{\textbf{Домашняя работа 3 по дискретной математике}\\[0.5em]
\large Нахождение пути максимальной пропускной способности.\\
Алгоритм Франка\,--\,Фриша (7.1.5)}
\author{Кливцов Александр Александрович \\ ISU: 504590, Вариант 29}
\date{}

\begin{document}

\maketitle
\thispagestyle{empty}

% ──────────────────────────────────────────────
\section*{Условие задачи}

Для взвешенного графа $G$ с вершинами $e_1, \ldots, e_{12}$, заданного матрицей
смежности~$R$, найти путь максимальной пропускной способности (максимального
минимального веса рёбер) от вершины $s = e_1$ до вершины $t = e_9$, используя
\textbf{алгоритм Франка\,--\,Фриша (7.1.5)}.

% ──────────────────────────────────────────────
\section*{Исходная матрица смежности $R$}

\[
	R =
	\renewcommand{\arraystretch}{1.1}
	\begin{array}{c|cccccccccccc}
		       & e_1 & e_2 & e_3 & e_4 & e_5 & e_6 & e_7 & e_8 & e_9 & e_{10} & e_{11} & e_{12} \\\hline
		e_1    & 0   & 4   & 4   &     & 2   &     &     & 2   &     &        & 3      &        \\
		e_2    & 4   & 0   &     &     & 4   &     & 2   &     & 5   &        & 2      & 2      \\
		e_3    & 4   &     & 0   &     &     & 4   & 4   & 4   &     & 5      & 2      &        \\
		e_4    &     &     &     & 0   &     & 5   &     &     & 4   &        & 4      &        \\
		e_5    & 2   & 4   &     &     & 0   & 3   &     &     &     &        &        & 4      \\
		e_6    &     &     & 4   & 5   & 3   & 0   &     & 4   &     & 3      & 2      &        \\
		e_7    &     & 2   & 4   &     &     &     & 0   &     &     & 1      & 3      &        \\
		e_8    & 2   &     & 4   &     &     & 4   &     & 0   &     & 1      & 3      &        \\
		e_9    &     & 5   &     & 4   &     &     &     &     & 0   &        &        &        \\
		e_{10} &     &     & 5   &     &     & 3   & 1   & 1   &     & 0      &        &        \\
		e_{11} & 3   & 2   & 2   & 4   &     & 2   & 3   & 3   &     &        & 0      &        \\
		e_{12} &     & 2   &     &     & 4   &     &     &     &     &        &        & 0      \\
	\end{array}
\]

% ──────────────────────────────────────────────
\section*{Список смежности (с весами)}

\begin{center}
	\begin{tabular}{ll}
		$\Gamma e_1 = \{e_2{:}4,\ e_3{:}4,\ e_5{:}2,\ e_8{:}2,\ e_{11}{:}3\}$              &
		$\Gamma e_7 = \{e_2{:}2,\ e_3{:}4,\ e_{10}{:}1,\ e_{11}{:}3\}$                            \\
		$\Gamma e_2 = \{e_1{:}4,\ e_5{:}4,\ e_7{:}2,\ e_9{:}5,\ e_{11}{:}2,\ e_{12}{:}2\}$ &
		$\Gamma e_8 = \{e_1{:}2,\ e_3{:}4,\ e_6{:}4,\ e_{10}{:}1,\ e_{11}{:}3\}$                  \\
		$\Gamma e_3 = \{e_1{:}4,\ e_6{:}4,\ e_7{:}4,\ e_8{:}4,\ e_{10}{:}5,\ e_{11}{:}2\}$ &
		$\Gamma e_9 = \{e_2{:}5,\ e_4{:}4\}$                                                      \\
		$\Gamma e_4 = \{e_6{:}5,\ e_9{:}4,\ e_{11}{:}4\}$                                  &
		$\Gamma e_{10} = \{e_3{:}5,\ e_6{:}3,\ e_7{:}1,\ e_8{:}1\}$                               \\
		$\Gamma e_5 = \{e_1{:}2,\ e_2{:}4,\ e_6{:}3,\ e_{12}{:}4\}$                        &
		$\Gamma e_{11} = \{e_1{:}3,\ e_2{:}2,\ e_3{:}2,\ e_4{:}4,\ e_6{:}2,\ e_7{:}3,\ e_8{:}3\}$ \\
		$\Gamma e_6 = \{e_3{:}4,\ e_4{:}5,\ e_5{:}3,\ e_8{:}4,\ e_{10}{:}3,\ e_{11}{:}2\}$ &
		$\Gamma e_{12} = \{e_2{:}2,\ e_5{:}4\}$                                                   \\
	\end{tabular}
\end{center}

% ──────────────────────────────────────────────
\section*{Выполнение алгоритма Франка\,--\,Фриша}

Напомним: на каждом шаге находим \textbf{максимальный} вес рёбер разреза $Q_k$,
затем \textbf{закорачиваем} все рёбра разреза с весом $\geq Q_k$
(объединяем их концевые вершины в одну супер-вершину).
Алгоритм завершается, когда $s$ и $t$ окажутся в одной супер-вершине.
Пропускная способность оптимального пути равна $\min(Q_1, Q_2, \ldots)$.

% ══════════════════════
\subsection*{Шаг 1. Разрез $K_1 = \bigl(\{e_1\},\ X \setminus \{e_1\}\bigr)$}

\textbf{Рёбра разреза} (рёбра, выходящие из $e_1$):

\begin{center}
	\begin{tabular}{|c|c|c|c|c|}
		\hline
		Ребро        & $(e_1,e_2)$ & $(e_1,e_3)$ & $(e_1,e_5)$ & $(e_1,e_8)$ \\
		\hline
		Вес $q_{ij}$ & 4           & 4           & 2           & 2           \\
		\hline
	\end{tabular}
	\quad
	\begin{tabular}{|c|c|}
		\hline
		Ребро        & $(e_1,e_{11})$ \\
		\hline
		Вес $q_{ij}$ & 3              \\
		\hline
	\end{tabular}
\end{center}

\[
	Q_1 = \max_{(x_i,x_j)\in K_1}[q_{ij}] = \max(4,\,4,\,2,\,2,\,3) = \mathbf{4}
\]

\textbf{Закорачиваем} все рёбра разреза $K_1$ с весом $\geq Q_1 = 4$:
$(e_1,e_2) = 4$ и $(e_1,e_3) = 4$.

Объединяем вершины:
\[
	v_1 = \{e_1,\,e_2,\,e_3\}, \quad
	\Gamma v_1 = \Gamma e_1 \cup \Gamma e_2 \cup \Gamma e_3 \;\setminus\; \{e_1,e_2,e_3\}.
\]

\textbf{Веса рёбер из супер-вершины} $v_1$ в $G_1$ (берём максимум):

\begin{center}
	\renewcommand{\arraystretch}{1.2}
	\begin{tabular}{|l|c|l|}
		\hline
		\textbf{Сосед} & \textbf{Вес $v_1$–сосед}                                             & \textbf{Исходное ребро (с макс. весом)} \\
		\hline
		$e_5$          & $\max(e_1{-}e_5{:}2,\; e_2{-}e_5{:}4) = 4$                           & $e_2 - e_5$                             \\
		$e_6$          & $\max(e_3{-}e_6{:}4) = 4$                                            & $e_3 - e_6$                             \\
		$e_7$          & $\max(e_2{-}e_7{:}2,\; e_3{-}e_7{:}4) = 4$                           & $e_3 - e_7$                             \\
		$e_8$          & $\max(e_1{-}e_8{:}2,\; e_3{-}e_8{:}4) = 4$                           & $e_3 - e_8$                             \\
		$e_9$          & $\max(e_2{-}e_9{:}5) = \mathbf{5}$                                   & $e_2 - e_9$                             \\
		$e_{10}$       & $\max(e_3{-}e_{10}{:}5) = \mathbf{5}$                                & $e_3 - e_{10}$                          \\
		$e_{11}$       & $\max(e_1{-}e_{11}{:}3,\; e_2{-}e_{11}{:}2,\; e_3{-}e_{11}{:}2) = 3$ & $e_1 - e_{11}$                          \\
		$e_{12}$       & $\max(e_2{-}e_{12}{:}2) = 2$                                         & $e_2 - e_{12}$                          \\
		\hline
	\end{tabular}
\end{center}

Граф $G_1$ имеет 10 вершин: $v_1, e_4, e_5, e_6, e_7, e_8, e_9, e_{10}, e_{11}, e_{12}$.
Вершины $s = e_1 \in v_1$ и $t = e_9 \notin v_1$ --- разделены, продолжаем.

% ══════════════════════
\subsection*{Шаг 2. Разрез $K_2 = \bigl(v_1,\ V(G_1) \setminus v_1\bigr)$ в графе $G_1$}

\textbf{Рёбра разреза} $K_2$ (рёбра, выходящие из $v_1$):

\begin{center}
	\renewcommand{\arraystretch}{1.2}
	\begin{tabular}{|c|c|c|c|c|c|c|c|c|}
		\hline
		Ребро & $v_1{-}e_5$ & $v_1{-}e_6$ & $v_1{-}e_7$ & $v_1{-}e_8$ & $v_1{-}e_9$ & $v_1{-}e_{10}$ & $v_1{-}e_{11}$ & $v_1{-}e_{12}$ \\
		\hline
		Вес   & 4           & 4           & 4           & 4           & \textbf{5}  & \textbf{5}     & 3              & 2              \\
		\hline
	\end{tabular}
\end{center}

\[
	Q_2 = \max_{(x_i,x_j)\in K_2}[q_{ij}] = \max(4,4,4,4,5,5,3,2) = \mathbf{5}
\]

\textbf{Закорачиваем} все рёбра разреза $K_2$ с весом $\geq Q_2 = 5$:
\begin{itemize}
	\item $v_1 - e_9 = 5$ (исходное ребро $e_2{-}e_9$): объединяем $e_9$ с $v_1$.
	\item $v_1 - e_{10} = 5$ (исходное ребро $e_3{-}e_{10}$): объединяем $e_{10}$ с $v_1$.
\end{itemize}

\[
	v_2 = v_1 \cup \{e_9,\,e_{10}\} = \{e_1,\,e_2,\,e_3,\,e_9,\,e_{10}\}.
\]

Проверяем: $s = e_1 \in v_2$ и $t = e_9 \in v_2$ --- \textbf{вершины $s$ и $t$ объединены. Алгоритм завершён.}

% ──────────────────────────────────────────────
\section*{Нахождение оптимального пути}

Пропускная способность оптимального пути:
\[
	\boxed{Q^* = \min(Q_1,\, Q_2) = \min(4,\, 5) = 4}
\]

\textbf{Граф $G'$} образован вершинами исходного графа и закороченными рёбрами:

\begin{center}
	\begin{tabular}{|c|c|c|}
		\hline
		\textbf{Шаг} & \textbf{Закороченные рёбра} & \textbf{Вес} \\
		\hline
		1            & $e_1 - e_2$                 & 4            \\
		1            & $e_1 - e_3$                 & 4            \\
		2            & $e_2 - e_9$                 & 5            \\
		2            & $e_3 - e_{10}$              & 5            \\
		\hline
	\end{tabular}
\end{center}

Ищем $(s$-$t)$ путь в $G'$ от $e_1$ до $e_9$:

\[
	e_1 \xrightarrow{4} e_2 \xrightarrow{5} e_9
\]

\[
	\text{Пропускная способность пути} = \min(4,\, 5) = \mathbf{4} \quad \checkmark
\]

% ──────────────────────────────────────────────
\section*{Проверка оптимальности}

Все пути из $e_1$ в $e_9$ в исходном графе:

\begin{center}
	\renewcommand{\arraystretch}{1.3}
	\begin{tabular}{|l|c|}
		\hline
		\textbf{Путь}                                 & \textbf{Пропускная способность} \\
		\hline
		$e_1 \to e_2 \to e_9$                         & $\min(4,5) = \mathbf{4}$        \\
		$e_1 \to e_3 \to e_6 \to e_4 \to e_9$         & $\min(4,4,5,4) = 4$             \\
		$e_1 \to e_{11} \to e_4 \to e_9$              & $\min(3,4,4) = 3$               \\
		$e_1 \to e_2 \to e_{11} \to e_4 \to e_9$      & $\min(4,2,4,4) = 2$             \\
		$e_1 \to e_8 \to e_3 \to e_6 \to e_4 \to e_9$ & $\min(2,4,4,5,4) = 2$           \\
		\hline
	\end{tabular}
\end{center}

Максимальная пропускная способность равна $\mathbf{4}$, что достигается на пути
$e_1 \to e_2 \to e_9$ (и также на $e_1 \to e_3 \to e_6 \to e_4 \to e_9$). Результат совпадает с ответом алгоритма. $\checkmark$

\bigskip
\noindent\textbf{Ответ:} оптимальный путь максимальной пропускной способности ---
$e_1 \to e_2 \to e_9$, его пропускная способность равна $\boldsymbol{4}$.

\end{document}
